% 示例阅读顺序：查看实际效果，在源码中复制对应 frame，再按需替换内容。
\section{分栏与图片}
\subsection{分栏排版}
\begin{frame}{双栏与不等宽分栏}
  \begin{columns}[T,onlytextwidth]
    \column{0.56\textwidth}
      \textbf{较宽的文字栏}
      \begin{enumerate}
        \item 使用 columns 创建分栏。
        \item 用 column 指定每栏宽度。
        \item 两栏之间预留适当空隙。
      \end{enumerate}
    \column{0.40\textwidth}
      \begin{block}{较窄的内容栏}
        放置结论、公式或图片，顶部与左栏对齐。
      \end{block}
  \end{columns}
  \slideNote{使用 [T,onlytextwidth] 顶部对齐并限制总宽度；
  栏宽之和小于正文宽度，为栏间留出空隙。}
\end{frame}

\subsection{插图与图题}
% 用法：\slideimage[高度上限]{文件名}；来源说明放在对应图片所在栏内。
\begin{frame}{插图与来源标注}
  \begin{columns}[T,onlytextwidth]
    \column{0.54\textwidth}
      \centering
      \slideimage[0.5\textheight]{example-figure.png}
      \sourceNote{来源：模板内置占位图。}
    \column{0.42\textwidth}
      \begin{itemize}
        \item \texttt{\textbackslash slideimage\{文件名\}}：按栏宽等比缩放。
        \item 加可选参数可限制高度，例如 \texttt{[0.45\textbackslash textheight]}。
        \item 图片统一放在 \texttt{figures/}，只写文件名即可。
      \end{itemize}
  \end{columns}
\end{frame}

\begin{frame}{并排两图}
  \begin{columns}[T,onlytextwidth]
    \column{0.48\textwidth}
      \centering
      \begin{figure}
        \slideimage[0.36\textheight]{example-figure.png}
        \caption{方案 A}
      \end{figure}
    \column{0.48\textwidth}
      \centering
      \begin{figure}
        \slideimage[0.36\textheight]{example-figure.png}
        \caption{方案 B}
      \end{figure}
  \end{columns}
  \sourceNote{来源：同一占位图重复使用，正式汇报时替换为两组真实结果。}
\end{frame}

% caption 之后添加 label；通过 ref 引用实际编号，不手写图号。
\begin{frame}{带图题与交叉引用的图片}
  \begin{figure}
    \centering
    \slideimage[0.36\textheight]{example-figure.png}
    \caption{示例图片与图题}\label{fig:demo}
    \sourceNote{来源：模板内置占位图。}
  \end{figure}
  图~\ref{fig:demo} 使用 figure、caption 和 label 添加图题与编号引用。
\end{frame}
